\section{Introduction}

TPC-65 reduced the represented-reference gate to the least eigenvalue of an
actual activated-range Gram \citep{WangTPC65}. TPC-66 then opened its literal
fixed-shift geometry: after retaining every active physical row and every full
output key \((\gamma,s,j)\), the normalized row columns have norm one; fixed
row pairs meet in only \(X^{o(1)}\) reduced cells on the critical cofactor
block \citep{WangTPC66}. That is genuine pair thinning, but not a lower-frame
theorem. The full \(\gamma\)-decorated support, right degrees, energy
concentration, and global cycles can still be large.

This paper asks the smallest next question:
\begin{quote}
\emph{Can the global lower frame be certified by the determinant mass of the
actual full-output matrix, and exactly which collision geometry controls that
mass?}
\end{quote}
The answer has a positive exact part and a sharp negative boundary.

First, a square physical minor is reduced to its nonpeelable collision core.
Every degree-one support vertex forces one matching pivot and one Laplace
factor. Once a perfect matching of the core is placed on the diagonal, the
normalized matrix is \(I+E\). The determinant expansion is then not an
unstructured sum: every nonidentity term is exactly a vertex-disjoint family
of directed alternating cycles. A diagonal similarity removes all phases on a
spanning forest, leaving only chord holonomies. This generalizes the
unique-matching and one-cycle examples of TPC-66 without replacing the literal
coefficients by generic fillings.

Second, the absolute activity
\[
 \mathcal A_\phi=\operatorname{per}(I+|E|)-1
\]
gives a sharp dominant-matching certificate. If
\(\mathcal A_\phi<1\), the chosen matching term cannot be canceled by all
cycle families combined. A useful sufficient version sums simple directed
cycle weights. Its \(\log 2\) threshold is deliberately presented only as a
consequence of an exponential majorant, not as a sharp physical threshold.

Third, we avoid the main quantitative defect of a one-minor argument.
If \(B\) has \(r\) unit columns and \(K=B^*B\), all full-column minors satisfy
\[
 \sum_{|Y_0|=r}|\det B_{Y_0}|^2=\det K.
\]
Since \(\operatorname{tr}K=r\), the other \(r-1\) eigenvalues can consume only
the remaining trace. For \(r\ge2\) this yields
\[
 \lambda_{\min}(K)
 \ge
 c_r\det K,
 \qquad
 c_r=\left(\frac{r-1}{r}\right)^{r-1}\ge e^{-1}.
\]
while \(c_1=1\). The constant is sharp. Hence a fixed-power lower bound for the aggregate
determinant mass transfers with zero exponent loss, unlike the
\(\|T\|^{-2(r-1)}\) conversion attached to one arbitrary minor.

The limitation is equally important. Determinant mass multiplies all
eigenvalues. It may therefore be exponentially small in the number of rows
even when the least eigenvalue stays uniformly positive. We exhibit this
behavior on a path Gram. Failure of the determinant route is consequently not
a physical spectral obstruction; it is a certified reason to return to the
direct comparison/forest route of TPC-66 or to open the later signed
reassembly door.

\subsection*{Claim levels}
The peeling, alternating-cycle expansion, holonomy count, activity bounds,
Cauchy--Binet identity, trace conversion, and countermodels are L0. Identifying
their matrix with the literal fixed-\(h_0\) full-output synthesis and retaining
all active rows is L1. A polynomial lower bound for the actual determinant
mass or physical Gram gap would be L2 and is not proved. Reduced cells are not
substituted for \((\gamma,s,j)\), no random phase or generic filling is used,
and no property special to \(h_0=2\) enters.
